\documentclass[11pt]{article}

\usepackage[a4paper,margin=2.5cm]{geometry}
\usepackage{amsmath}
\usepackage{amssymb}
\usepackage{xcolor}

\newcommand{\CS}{\textcolor{red}{\mathbf{CS}}}
\newcommand{\Score}{\textcolor{blue}{\mathrm{Score}}}
\newcommand{\MLM}{\textcolor{blue}{\mathcal{M}}}

\title{Construction Explained}
\author{Kirill Shmakov}
\date{\today}

\begin{document}

\nocite{*}

\maketitle

%%%%%%%%%%%%%%%%%%%%%%%%%%%%%%%%%%%%%%%%%%%%%%%%%%%%%%%%%%%%%%%%%%%%%%%%%%%%%%%%

\section{Goal}

The idea of this project is to augment the algorithm presented
in~\cite{rokach2013mlcircuit} with an ML-backed model for complexity estimation
of decision trees, used for the heuristic choice of element.

\subsection{Formulation}

We work over a fixed set of logical elements
$\mathcal{K} = \left\{ K_1, K_2, \dots, K_m \right\}$, each of which is a Boolean map
$$
    K : \left\{0, 1\right\}^{\mathrm{in}(K)} \to \left\{0, 1\right\}^{\mathrm{out}(K)}.
$$
We are given a Boolean function $g$ with $n$ primary inputs $\bar{x} = \left(x_1, x_2, \dots, x_n\right)$,
and our goal is to construct a circuit representation of $g$ using elements from $\mathcal{K}$.

%%%%%%%%%%%%%%%%%%%%%%%%%%%%%%%%%%%%%%%%%%%%%%%%%%%%%%%%%%%%%%%%%%%%%%%%%%%%%%%%

\section{Algorithm}

\subsection{Baseline}

The construction of~\cite{rokach2013mlcircuit}, there called the Circuit Decomposition
Engine, is a greedy heuristic that grows the circuit one element at a time, moving from
the primary inputs towards the output. The framework of this algorithm is as follows.

\paragraph{State.} The algorithm operates on a pair $(S, C)$, where $S$ is the set of
\emph{slots} -- the signals that are available -- and $C$ is the current unfinished
circuit. Every slot $y \in S$ is itself a Boolean function of the primary inputs,
so $S$ can be read as a value table with one column per slot. Initially
$$
    S_0 = \left\{x_1, x_2, \dots, x_n\right\}, \qquad C_0 = \text{the empty circuit}.
$$

\paragraph{Residual.} Together with $g$, a slot set $S$ determines the \emph{residual
problem}: computing $g$ from the columns of $S$ alone. The score function
$\Score\left[g \mid S\right]$ measures the complexity of that residual, and the greedy
step is the attempt to make it drop as fast as possible.

\paragraph{Step.} For the current pair $(S_j, C_j)$ the algorithm enumerates every candidate
attachment: an element $K \in \mathcal{K}$ with $k = \mathrm{in}(K)$, together with a
$k$-subset $\left\{y_{i_1}, \dots, y_{i_k}\right\} \subseteq S$ of slots to feed it. The
element contributes all of its $l = \mathrm{out}(K)$ outputs at once, as $l$ new slots
$$
    \left(z_1, \dots, z_l\right) = K\left(y_{i_1}, \dots, y_{i_k}\right),
$$
whose columns are computed row-wise from the columns of the inputs, and the candidate
carries the state to
$$
    S_{j+1} = S_j \cup \left\{z_1, \dots, z_l\right\},
    \qquad
    C_{j+1} = C_j + K.
$$
Every candidate is rated by $\Score\left[g \mid S_{j+1}\right]$, the attachment with the
minimal value is committed, and the algorithm proceeds from $(S_{j+1}, C_{j+1})$.

\paragraph{Termination.} The algorithm stops once one of the slots computes $g$.

\subsection{Limitation}

In the original approach $\Score\left[g \mid S\right]$ is the number of leaves in a
decision tree that computes $g$ from the value table of $S$. Finding a tree with the
smallest number of leaves is NP-complete~\cite{hyafil1976np}, so computing this score
exactly costs time exponential in the arity.

To speed this up, the authors used Orange's \texttt{TreeLearner} heuristic, which returns
an approximate smallest tree. Even so, the heuristic still has to process every row of
the value table, and the number of rows grows exponentially with the arity of $g$. The
choice gives good convergence, but it quickly becomes impractical as that arity grows.

Formally, the tree has to be built from a table of size $2^n \times L$, where $n$ is the
arity of $g$ and $L$ is the number of slots, itself proportional to the size of the current
circuit $C$. Both dimensions grow exponentially in $n$: the rows by definition, and the
columns because the classical results of Shannon~\cite{shannon1949synthesis} and
Lupanov~\cite{lupanov1958synthesis} put the circuit complexity of almost all Boolean
functions of $n$ variables at order $2^n / n$.

\subsection{Our modification}

Instead of working with the exact decision tree, we take the score to be the
approximate value
$$
    \Score\left[g \mid S\right] = \MLM\left[g \mid S\right],
$$
backed by a learned function $\MLM$. The main benefit of this approach is that $\Score$
becomes computationally affordable, which lets us handle $g$ of larger arity.

\section{Model}

For $n \leqslant 12$ we can compute $\CS\left[g \mid S\right]$ as the size
of the minimal decision tree (the truth table is not very big, so dynamic programming
is applicable).

For larger $n$ this quickly becomes impossible, and we need to use an ML model $S_1$
for approximate computation.

$S_1$ accepts two boolean functions $g, f : \{0,1\}^n \to \{0,1\}$ as its arguments;
we write $S^n_1$ when we need to track the arity of the input functions.
It outputs a pair $(d, s)$ of real numbers, $d, s \in [0, n]$. Here $d$ is the depth
of the minimal-depth decision tree, and $s$ is $\log_2 (\text{IN} + 1)$, where
$\text{IN}$ is the number of internal nodes in the minimal-size decision tree; these
two trees need not be the same.

We also need another model $S_2$, which takes a boolean function $g : \{0,1\}^n \to \{0,1\}$
and a subset $X \subseteq \{0,1\}^n$ as its arguments. Its outputs $(d, s)$ have the
same semantics, for the decision trees needed to compute $g$ restricted to $X$.

For training $S_1$ and $S_2$ we use exact values when possible. Otherwise we compute
the target value as follows.

\subsection{Targets for $S_1$}

We select the target for $S^n_1\left[g \mid f\right]$ to be the minimum over the
following reductions:
\begin{itemize}
    \item through each primary input $x_i$:
        $$
            S^{n-1}_1 \left[g \downarrow_{x_i = 0} \mid f \downarrow_{x_i = 0} \right], \qquad
            S^{n-1}_1 \left[g \downarrow_{x_i = 1} \mid f \downarrow_{x_i = 1} \right]
        $$
    \item through $f$:
        $$
            S^n_2 \left[g \mid f^{-1}(0) \right], \qquad
            S^n_2 \left[g \mid f^{-1}(1) \right]
        $$
\end{itemize}

Here $h \downarrow_{x_i = \alpha}$ for $\alpha \in \{0, 1\}$ denotes the restriction of
$h$ obtained by fixing $x_i = \alpha$, viewed as a function of the remaining $n - 1$
inputs.

\subsection{Targets for $S_2$}

We select the target for $S^n_2\left[g \mid X\right]$ to be the minimum over the
reductions through each primary input $x_i$:
$$
    S^{n-1}_2 \left[g \downarrow_{x_i = 0} \mid X \downarrow_{x_i = 0} \right], \qquad
    S^{n-1}_2 \left[g \downarrow_{x_i = 1} \mid X \downarrow_{x_i = 1} \right]
$$
where $X \downarrow_{x_i = \alpha}$ is the projection of $X \cap \{x_i = \alpha\}$ onto
the coordinates other than $x_i$.

\bibliographystyle{unsrt}
\bibliography{paper}

\end{document}
